% -*- coding: utf-8 -*-
% !TEX program = xelatex
\documentclass[plain,noanswer,medmath]{../../template/jnuexam}
\usepackage{../../template/kysx}

\begin{document}


\examtitle{1990}{4}


\makepart{填空题}{本题满分15分，每小题3分}

\begin{problem}
极限$\lim_{n\to\infty}\,\left(\sqrt{n+3\sqrt{n}}-\sqrt{n-\sqrt{n}}\right)=$\fillin{}.
\end{problem}


\begin{problem}
设函数$f(x)$有连续的导函数，$f(0)=0,f'(0)=b$，若函数
$$F(x)=\begin{cases}
  \frac{f(x)+a\sin x}{x}, & x\ne 0, \\
                       A, & x=0
\end{cases}$$
在$x=0$处连续，则常数$A$=\fillin{}.
\end{problem}


\begin{problem}
曲线$y=x^2$与直线$y=x+2$所围成的平面图形的面积为\fillin{}.
\end{problem}


\begin{problem}
若线性方程组$\begin{cases}
  x_1+x_2=-a_1, \\
  x_2+x_3=a_2, \\
  x_3+x_4=-a_3, \\
  x_4+x_1=a_4 \\
\end{cases}$有解，则常数$a_1$, $a_2$, $a_3$, $a_4$应满足条件\fillin{}.
\end{problem}


\begin{problem}
一射手对同一目标独立地进行四次射击，若至少命中一次的概率为$\frac{80}{81}$,
则该射手的命中率为\fillin{}.
\end{problem}


\makepart{选择题}{本题满分15分，每小题3分}


\begin{problem}
设函数$f(x)=x\cdot \tan x\cdot \e^{\sin x}$，则$f(x)$是\pickout{}
\begin{abcd}
\item 偶函数．
\item 无界函数．
\item 周期函数．
\item 单调函数．
\end{abcd}
\end{problem}


\begin{problem}
设函数$f(x)$对任意$x$均满足等式$f(1+x)=af(x)$，且有$f'(0)=b$,
其中$a$, $b$为非零常数，则\pickout{}
\begin{abcd}
\item $f(x)$在$x=1$处不可导．
\item $f(x)$在$x=1$处可导，且$f'(1)=a$．
\item $f(x)$在$x=1$处可导，且$f'(1)=b$．
\item $f(x)$在$x=1$处可导，且$f'(1)=ab$．
\end{abcd}
\end{problem}


\begin{problem}
向量组$\alpha_1,\alpha_2,\cdots,\alpha_s$线性无关的充分条件是\pickout{}
\begin{abcd}
\item $\alpha_1,\alpha_2,\cdots,\alpha_s$均不为零向量．
\item $\alpha_1,\alpha_2,\cdots,\alpha_s$中任意两个向量的分量不成比例．
\item $\alpha_1,\alpha_2,\cdots,\alpha_s$中任意一个向量均不能由其余$s-1$个向量线性表示．
\item $\alpha_1,\alpha_2,\cdots,\alpha_s$中有一部分向量线性无关．
\end{abcd}
\end{problem}


\begin{problem}
设$A,B$为两随机事件，且$B\subset A$，则下列式子正确的是\pickout{}
\begin{abcd}
\item $P(A+B)=P(A)$．
\item $P(AB)=P(A)$．
\item $P(B|A)=P(B)$．
\item $P(B-A)=P(B)-P(A)$．
\end{abcd}
\end{problem}


\begin{problem}
设随机变量$X$和$Y$相互独立，其概率分布为
$$\begin{array}{|c|cc|}
\hline
         m &  -1 & 1 \\
\hline
  P\{X=m\} & 1/2 & 1/2\\
\hline
\end{array}
\qquad
\begin{array}{|c|cc|}
\hline
         m &  -1 & 1 \\
\hline
  P\{Y=m\} & 1/2 & 1/2 \\
\hline
\end{array}$$
则下列式子正确的是\pickout{}
\begin{abcd}
\item $X=Y$．
\item $P\{X=Y\}=0$．
\item $P\{X=Y\}=\frac{1}{2}$．
\item $P\{X=Y\}=1$．
\end{abcd}
\end{problem}


\makepart{计算题}{本题满分20分，每小题5分}


\begin{problem}
求函数$I(x)=\int_{\e}^x \frac{\ln t}{t^2-2t+1}\dt$在区间$[\e, \e^2]$上的最大值．
\end{problem}


\begin{problem}
计算二重积分$\iint_D x\e^{-y^2}\dx\dy$，其中$D$是曲线$y=4x^2$和$y=9x^2$在第一象限所围成的区域．
\end{problem}


\begin{problem}
求级数$\sum_{n=1}^{\infty}\frac{(x-3)^n}{n^2}$的收敛域．
\end{problem}


\begin{problem}
求微分方程$y'+y\cos x=(\ln x)\e^{-\sin x}$的通解．
\end{problem}


\makepart{}{本题满分9分}


\begin{problem}
某公司可通过电台及报纸两种形式做销售某种商品的广告，根据统计资料，
销售收入$R$(万元)与电台广告费用$x_1$(万元)及报纸广告费用$x_2$(万元)之间的关系有如下经验公式:
$$R=15+14x_1+32x_2-8x_1x_2-2x_1^2-10x_2^2.$$
\begin{enumerate}
\item 在广告费用不限的情况下，求最优广告策略；
\item 若提供的广告费用为$1.5$万元，求相应的最优广告策略．
\end{enumerate}
\end{problem}


\makepart{}{本题满分6分}


\begin{problem}
设$f(x)$在闭区间$[0,c]$上连续，其导数$f'(x)$在开区间$(0,c)$内存在且单调减少；$f(0)=0$，
试应用拉格朗日中值定理证明不等式：$f(a+b)\le f(a)+f(b)$，其中常数$ab$满足条件$0\le a\le b\le a+b\le c$．
\end{problem}


\makepart{}{本题满分8分}


\begin{problem}
已知线性方程组$\begin{cases}
  x_1+x_2+x_3+x_4+x_5=a, \\
  3x_1+2x_2+x_3+x_4-3x_5=0, \\
  x_2+2x_3+2x_4+6x_5=b, \\
  5x_1+4x_2+3x_3+3x_4-x_5=2,
\end{cases}$
\begin{enumerate}
\item $a$, $b$为何值时，方程组有解？
\item 方程组有解时，求出方程组的导出组的一个基础解系；
\item 方程组有解时，求出方程组的全部解．
\end{enumerate}
\end{problem}


\makepart{}{本题满分5分}

\begin{problem}
已知对于$n$阶方阵$A$，存在自然数$k$，使得$A^k=0$，试证明矩阵$E-A$可逆，
并写出其逆矩阵的表达式($E$为$n$阶单位阵).
\end{problem}


\makepart{}{本题满分6分}


\begin{problem}
设$A$是$n$阶矩阵，$\lambda_1$和$\lambda_2$是$A$的两个不同的特征值，
$x_1,x_2$是分别属于$\lambda_1$和$\lambda_2$的特征向量．
试证明$x_1+x_2$不是$A$的特征向量．
\end{problem}


\makepart{}{本题满分4分}


\begin{problem}
从$0,1,2,\cdots,9$十个数字中任意选出三个不同数字，试求下列事件的概率：\par
\qquad$A_1=$\{三个数字中不含$0$和$5$\}；\quad$A_2=$\{三个数字中不含$0$或$5$\}.
\end{problem}


\makepart{}{本题满分5分}


\begin{problem}
一电子仪器由两个部件构成，以$X$和$Y$分别表示两个部件的寿命（单位：千小时），
已知$X$和$Y$的联合分布函数为：
$$F(x,y)=\begin{cases}
  1-\e^{-0.5x}-\e^{-0.5y}+\e^{-0.5(x+y)}, & x\ge 0, y\ge 0, \\
                                       0, & \text{其他}. \\
\end{cases}$$
\begin{enumerate}
\item 问$X$和$Y$是否独立？
\item 求两个部件的寿命都超过100小时的概率$\alpha$．
\end{enumerate}
\end{problem}


\makepart{}{本题满分7分}


\begin{problem}
某地抽样调查结果表明，考生的外语成绩(百分制)近似服从正态分布，平均成绩为72分，96分以上的占考生总数的2.3\%,
试求考生的外语成绩在60分至84分之间的概率．\par
\centerline{[附表]（表中$\Phi (x)$是标准正态分布函数．）}\noindent
$$\begin{array}{|c|ccccccc|}
\hline
        x &     0 &   0.5 &   1.0 &   1.5 &   2.0 &   2.5 & 3.0 \\
\hline
  \Phi(x) & 0.500 & 0.692 & 0.841 & 0.933 & 0.977 & 0.994 & 0.999 \\
\hline
\end{array}$$
\end{problem}


\end{document}
